%%%
%
% $Autor: Siddhanth Sharma (Matriculation No. 7027189) $
% $Datum: 2021-05-14 $
% $Dateiname: 
% $Version: 4620 $
%
% !TeX spellcheck = GB
% !TeX program = pdflatex
% !TeX encoding = utf8
%
%%%

\chapter{\textcolor{red}{PyTorch}}

\index{PyTorch}

% --------------------------------------------------------
\section{Introduction}
\index{PyTorch!Introduction}

\noindent
PyTorch is a modern, open-source deep learning framework designed to provide an intuitive, Pythonic interface while maintaining high computational performance. As described by its core contributors, PyTorch implements an imperative (define-by-run) execution model that enables flexible neural network construction and debugging during runtime \cite{paszke2019pytorch}. At its core, PyTorch offers a powerful tensor library comparable to NumPy but augmented with GPU acceleration and automatic differentiation through its \texttt{autograd} engine, which efficiently computes gradients for optimization in deep learning workloads \cite{pytorchdocs}.

What distinguishes PyTorch from earlier deep learning frameworks is its dynamic computation graph design, which constructs the computational graph during execution and allows developers to utilize standard Python control flow for model definition \cite{goodfellow2016deep}. Beyond its fundamental tensor and gradient capabilities, PyTorch provides a comprehensive ecosystem of modules for neural network layers, loss functions, optimization algorithms, and efficient data pipelines. These components, along with strong community support and extensive official documentation, make PyTorch a versatile and scalable framework for both research-oriented experimentation and large-scale production systems \cite{pinto2022springer}.



In this project, PyTorch serves as the primary deep learning framework responsible for executing the licence plate detection model on the Jetson Nano B01. After image frames are captured and pre-processed by OpenCV, PyTorch is used to:

\begin{itemize}
	\item Convert incoming frames into tensors suitable for neural network inference,
	\item Run the optimised deep learning model on the Jetson Nano’s CUDA-enabled GPU to achieve real-time licence plate detection, and
	\item Interface with \texttt{torchvision} utilities for model loading, tensor transformations, and streamlined inference workflows.
\end{itemize}

This tight coupling between PyTorch and the Jetson platform provides an efficient pipeline for identifying licence plates and logging detection results within the system \cite{jetsonpytorch}.



% --------------------------------------------------------
\section{Description}
\index{PyTorch!Description}

PyTorch is organised into a set of modular components that work together to support deep learning model development, training, and deployment. Each module provides specialised functionality, allowing developers to build complex pipelines with minimal overhead. The most relevant components include:

\begin{itemize}
	\item \texttt{torch}: the core package providing tensor operations, mathematical functions, and GPU acceleration.
	\item \texttt{torch.nn}: a high-level module containing neural network layers, activation functions, and loss functions.
	\item \texttt{torch.optim}: implements optimisation algorithms such as SGD and Adam for training neural networks.
	\item \texttt{torchvision}: offers tools for image-based tasks, including datasets, transforms, and pre-trained models.
	\item \texttt{torch.utils.data}: provides efficient data-loading utilities, including dataset abstractions and data loaders.
	\item \texttt{autograd}: PyTorch’s automatic differentiation engine that computes gradients for backpropagation.
\end{itemize}

\begin{figure}[htbp]
	\centering
	\begin{tikzpicture}[
		module/.style={
			rectangle,
			rounded corners,
			draw=black,
			thick,
			minimum width=3cm,
			minimum height=0.6cm,
			align=center,
			font=\small
		},
		core/.style={module, fill=blue!20},
		nn/.style={module, fill=green!20},
		optim/.style={module, fill=orange!20},
		util/.style={module, fill=yellow!20},
		arrow/.style={
			-{Latex[length=2.5mm]},
			thick
		}
		]
		
		% Main PyTorch box
		\node[core] (pytorch) {PyTorch Framework};
		
		% First level modules
		\node[core, below=1.5cm of pytorch, xshift=-5cm] (torch) {\texttt{torch}\\Core Tensor Library};
		\node[nn, below=1.5cm of pytorch] (nn) {\texttt{torch.nn}\\Neural Network Modules};
		\node[optim, below=1.5cm of pytorch, xshift=5cm] (optim) {\texttt{torch.optim}\\Optimisation Algorithms};
		
		% Second level modules under torch.nn, spread horizontally
		\node[util, below=2.0cm of nn, xshift=-5cm] (vision) {\texttt{torchvision}\\Image Tools \& Models};
		\node[util, below=2.0cm of nn] (data) {\texttt{torch.utils.data}\\Data Loading Utilities};
		\node[util, below=2.0cm of nn, xshift=5cm] (autograd) {\texttt{autograd}\\Automatic Differentiation};
		
		% Arrows from PyTorch to first level
		\draw[arrow] (pytorch.south west) -- (torch.north);
		\draw[arrow] (pytorch.south) -- (nn.north);
		\draw[arrow] (pytorch.south east) -- (optim.north);
		
		% Arrows from nn to second level
		\draw[arrow] (nn.south west) -- (vision.north);
		\draw[arrow] (nn.south) -- (data.north);
		\draw[arrow] (nn.south east) -- (autograd.north);
		
	\end{tikzpicture}
	\captionsetup{font=small}
	\caption{Hierarchical structure of the PyTorch framework with color-coded modules: core tensor operations (\texttt{torch}) in blue, neural network layers (\texttt{torch.nn}) in green, optimization algorithms (\texttt{torch.optim}) in orange, and supporting utilities (\texttt{torchvision}, \texttt{torch.utils.data}, \texttt{autograd}) in yellow.}
	\label{fig:pytorch_structure}
\end{figure}


PyTorch provides a highly flexible and efficient framework for computer vision applications. Its dynamic computation graph allows for rapid experimentation with custom neural network architectures, which is particularly valuable in vision tasks that often require complex, non-standard layers and pipelines \cite{paszke2019pytorch}. The seamless integration with Python and popular scientific libraries such as NumPy and SciPy further simplifies data manipulation, augmentation, and preprocessing, all of which are critical steps in building robust vision models \cite{goodfellow2016deep}.

PyTorch also offers a rich ecosystem of vision-specific utilities through its \texttt{torchvision} library. This module provides pre-built datasets, common image transforms, and pre-trained models such as ResNet, VGG, and YOLO, which can be easily fine-tuned for custom tasks \cite{torchvision_docs}. Additionally, PyTorch supports mixed-precision training, distributed computing, and optimized CUDA kernels, enabling efficient training and inference even on resource-constrained devices like edge GPUs and mobile platforms \cite{nvidia_jetson}. These capabilities allow developers to build state-of-the-art computer vision pipelines for tasks including object detection, image classification, and semantic segmentation with both speed and flexibility.

% --------------------------------------------------------
\section{Installation}
\index{PyTorch!Installation}

Installing PyTorch typically begins with preparing a standard Python environment. Users can install PyTorch using popular package managers such as \texttt{pip} or \texttt{conda}, which handle the library itself along with its Python bindings and essential dependencies. Core dependencies include NumPy, SciPy, and Python development headers, while GPU acceleration requires CUDA and cuDNN libraries on compatible hardware \cite{paszke2019pytorch, kulkarni2022cvprojects}. In a typical workflow, a user may first create a virtual environment using \texttt{venv} or \texttt{conda}, activate it, and then install PyTorch and optional packages like \texttt{torchvision}, \texttt{torchaudio}, or \texttt{torchtext} to build a complete deep learning stack. This approach ensures that PyTorch integrates seamlessly with other scientific computing libraries and that dependencies are isolated for project-specific workflows.

Installing PyTorch on an NVIDIA Jetson platform requires a series of concrete steps to ensure compatibility with the ARM architecture and the CUDA-enabled GPU. The typical installation procedure for Jetson Nano, TX2, or Xavier devices includes:

\begin{itemize}
	\item \textbf{Install NVIDIA JetPack SDK:} Download and install the JetPack SDK, which provides the necessary CUDA, cuDNN, and TensorRT libraries for GPU acceleration.
	\item \textbf{Set up Python environment:} Install Python 3 and \texttt{pip}, and optionally create a virtual environment using \texttt{virtualenv} or \texttt{conda} to isolate dependencies.
	\item \textbf{Install PyTorch wheel:} Download the pre-built PyTorch wheel specifically optimized for Jetson devices and install it using \texttt{pip}.
	\item \textbf{Install supporting libraries:} Install \texttt{torchvision} and other relevant packages ensuring CUDA support is enabled for GPU-accelerated operations.
	\item \textbf{Verify installation:} Run a simple tensor operation and confirm GPU availability to ensure PyTorch is fully operational.
\end{itemize}

Following these steps ensures that PyTorch is correctly installed and optimized for Jetson hardware, enabling efficient training and inference of deep learning models for computer vision, robotics, and other edge applications \cite{nvidia_jetson, kulkarni2022cvprojects}.

\begin{tcolorbox}[colback=white,colframe=black,coltext= ShellColor, sharp corners, boxrule=0.8pt]
	\begin{small}
		\begin{verbatim}
			
			pip install torch torchvision torchaudio	\end{verbatim}
	\end{small}
\end{tcolorbox}

This command installs the core \PYTHON{torch} library along with \PYTHON{torchvision} and \PYTHON{torchaudio} modules, which provide utilities for computer vision and audio processing tasks. After installation, these modules can be imported directly in Python scripts for model building, training, and inference. For advanced configurations, such as enabling CUDA support on desktop GPUs or Jetson devices, the official PyTorch documentation provides detailed installation instructions and platform-specific wheels \cite{paszke2019pytorch, pytorchdocs}.



% --------------------------------------------------------
\section{Example -- Description}
\index{OpenCV!Example Description}

The example script \FILE{PyTorchLPExample.py} demonstrates a minimal but representative PyTorch workflow for licence plate detection on the Jetson Nano. Its purpose is to verify that:
\begin{itemize}
	\item the camera is correctly connected and accessible,
	\item OpenCV can capture and display frames in real time,
	\item PyTorch tensors can be created from camera frames and moved to GPU if available, and
	\item basic model inference and frame annotation operations work as expected on the Jetson Nano.
\end{itemize}

Conceptually, the example performs the following steps:
\begin{enumerate}
	\item Open the default camera device using \PYTHON{cv2.VideoCapture(0)}.
	\item Read frames in a loop and check that the returned frame is valid.
	\item Convert each frame to a PyTorch tensor using \PYTHON{torchvision.transforms.ToTensor()}, and optionally move it to GPU with \PYTHON{tensor.cuda()}.
	\item Run a PyTorch model (dummy or pre-trained) on the tensor to simulate licence plate detection.
	\item Annotate the original frame with predicted bounding boxes and confidence scores using \PYTHON{cv2.rectangle} and \PYTHON{cv2.putText}.
	\item Display the annotated frame in a window using \PYTHON{cv2.imshow}.
	\item Terminate the loop when the user presses a specific key (e.g., \texttt{q}) and release all resources.
\end{enumerate}

This example mirrors the essential I/O, pre-processing, and inference steps used in the full Jetson Nano licence plate detection pipeline and provides a template for integrating PyTorch models with real-time camera input.


% --------------------------------------------------------
\section{Example -- Manual}
\index{OpenCV!Example Manual}

The following instructions describe how to run the PyTorch example on the Jetson Nano:

\begin{itemize}
	\item \textbf{Step 1:} Ensure that the Jetson Nano is powered on and that the camera (CSI or USB) is connected correctly as configured for the project.
	\item \textbf{Step 2:} Activate the Python virtual environment (if used) that contains the project dependencies, including PyTorch and torchvision.
	\item \textbf{Step 3:} Navigate to the project directory where \FILE{PyTorchLPExample.py} is stored, for example:
	
	\begin{tcolorbox}[colback=white, colframe=black, coltext=ShellColor, sharp corners, boxrule=0.8pt]
		\begin{small}
			\begin{verbatim}
				cd ~/projects/license-plate-detection/pytorch_tests
			\end{verbatim}
		\end{small}
	\end{tcolorbox}
	
	\item \textbf{Step 4:} Run the example script:
	
	\begin{tcolorbox}[colback=white, colframe=black, coltext=ShellColor, sharp corners, boxrule=0.8pt]
		\begin{small}
			\begin{verbatim}
				python PyTorchLPExample.py
			\end{verbatim}
		\end{small}
	\end{tcolorbox}
	
	\item \textbf{Step 5:} A window should open showing the camera stream with any model predictions (dummy bounding boxes or real detections) overlaid. Verify that the frame rate is smooth and that the bounding boxes are correctly displayed.
	\item \textbf{Step 6:} Press the configured exit key (e.g., \texttt{q}) to close the window, release the camera, and terminate the program.
\end{itemize}

This procedure provides a quick check of camera connectivity, PyTorch model inference, and real-time annotation before running the full Jetson Nano licence plate detection pipeline.


% --------------------------------------------------------
\section{Example -- Version}
\index{OpenCV!Example Version}

To ensure reproducibility and to document the exact PyTorch build used on the Jetson Nano, the following Python snippet can be executed. This script prints the version of the installed OpenCV package.

\subsection{Code}
\lstinputlisting[
language=Python,
caption={Checking the installed OpenCV version},
label={lst:opencv_version_check},
basicstyle=\ttfamily\small,
keywordstyle=\color{blue},
commentstyle=\color{gray},
stringstyle=\color{green!40!black},
breaklines=true,
showstringspaces=false,
numbers=left,
numberstyle=\tiny\color{gray},
frame=single,
rulecolor=\color{black!30}
]{../../Code/PyTorch/pytorchversion.py}

Recording the PyTorch version number in the project documentation helps to avoid compatibility issues when upgrading JetPack, CUDA, or when reinstalling packages on another device. For example, the online documentation pages used in this project correspond to PyTorch version 2.1.0, which provides the \emph{PyTorch Tutorials} and related examples for tensor operations, neural network building, and model inference \cite{pytorchdocs, paszke2019pytorch}.



% --------------------------------------------------------
\section{Example -- Files}
\index{OpenCV!Example Files}

The central example file for this section is:

\FILE{PyTorchLPExample.py}

This script is kept intentionally short and focused on:
\begin{itemize}
	\item initialising the camera,
	\item converting frames to PyTorch tensors and performing basic pre-processing, and
	\item visual verification of the annotated image stream with predicted bounding boxes.
\end{itemize}

In the final system, similar code fragments are integrated into the main application script that loads the YOLOv8 or custom PyTorch model, performs inference on each frame, and writes detection results to the SQLite database.


\subsection{Code}
\lstinputlisting[
language=Python,
caption={PyTorch Image Recognition Example for Jetson Nano},
label={lst:opencv_example},
basicstyle=\ttfamily\small,
keywordstyle=\color{blue},
commentstyle=\color{gray},
stringstyle=\color{green!40!black},
breaklines=true,
showstringspaces=false,
numbers=left,
numberstyle=\tiny\color{gray},
frame=single,
rulecolor=\color{black!30}
]{../../Code/PyTorch/pytorchexample.py}

% --------------------------------------------------------
\section{Further Reading}
\index{OpenCV!Further Reading}

The following references provide additional background and detailed documentation on OpenCV and its Python interface:

\nocite{Bradski:2000}
\nocite{opencv_wiki}
\nocite{opencv_docs}
\nocite{opencv_py_intro}
\nocite{opencv_py_tutorial_root}
\nocite{opencv_pypi}
\nocite{learning_opencv}

% Prints bibliography for this chapter
\printbibliography[heading=subbibliography, segment=\therefsegment]